\section*{Results}\label{section: result}
%% Overview of Renoir
\subsection*{Charting spatial ligand-target interactions using Renoir}
In order to build a comprehensive computation framework to investigate spatial cell-cell interaction we developed Renoir. As shown in Figure 1, Renoir is an end-to-end computational framework for exploring the spatial map of ligand-target activities either by integrating spatial transcriptomics and scRNA-seq datasets from the same tissue or by employing single-cell resolution spatial transcriptomics data. Renoir first curates a set of ligand-target pairs using NATMI's ConnectomeDB \citep{Hou2020} and NicheNet \citep{Browaeys2020}. For a low-resolution spatial transcriptomic dataset (e.g., 10X Visium), using a well-annotated scRNA-seq data, Renoir quantifies a neighborhood activity score for each curated ligand-target pair for each spot in the spatial transcriptomic data. For a single-cell resolution spatial transcriptomic data, Renoir quantifies the neighborhood activity score for each curated ligand-target pair for each cell with a spatial context. The quantification of the neighborhood activity score for a specific ligand-target pair for a particular spot is performed using the proportions and cell type-specific mRNA abundances present in each spot/cell within the defined neighborhood of the spot/cell (see Methods for details). The activity between a ligand-target pair for the given spot and another spot in the neighborhood is scored based on the cell type–specific mRNA abundances of the two genes weighted by the cell type proportions as well as the mutual information between the two genes across the cell types present within the two spots. Thus, Renoir generates a novel neighborhood activity score matrix (of dimensions number of ligand-target pairs $\times$ number of spots). Using the ligand-target neighborhood activity score matrix, Renoir performs several downstream analyses. First of all, Renoir infers spatial communication domains through unsupervised clustering of spots based on ligand-target neighborhood activity scores and determines communication domain-specific ligand-target activities. Moreover, Renoir can infer cell type-specific ligand-target activities within a communication domain. Renoir further provides an algorithm for ranking the ligands within a communication domain based on their cumulative activity across target genes expressed by the major cell types in the domain. Renoir also provides a mapping of pathway-level activity across spots and communication domains. Finally, Renoir encompasses a visualization module for the visualisation of ligand-target activity across spots as well as a Sankey plot that provides an overview of user-specified ligand-target activities across cell types and communication domains.

%% Renoir Benchmarking
\subsection*{Renoir outperforms state-of-the-art methods in inferring ligand-target interactions}
We first evaluated Renoir's performance in inferring the spatial activity of ligand-target pairs using simulated datasets generated using a semi-synthetic enrichment-based simulation workflow (Fig. 2a, see Methods for details). Specifically, we started with paired single-cell and spatial datasets based on which we simulated new spatial datasets where activity of a ligand-target pair was simulated in reference spatial locations (positive reference spots) through the enrichment of cell types associated with the ligand and target gene. Moreover, the simulated datasets also contained spatial locations which did not harbor any activity between the ligand and target (negative reference spots). To comprehensively cover the range of interactions across tissue types, we used three reference matched scRNA-seq and ST datasets from human intestine \citep{FAWKNERCORBETT2021810}, triple negative breast cancer (TNBC) \citep{Wu2021} and dorsolateral prefrontal cortex (DLPFC) \citep{Maynard2021} for the simulation. The simulation workflow as described in Figure 2a was applied to generate 43, 15 and 26 simulated datasets corresponding to the human intestine, TNBC and DLPFC datasets respectively, where, each simulated dataset enriched a specific ligand-target interaction. Since Renoir quantifies ligand-target activity at a specific spatial context and is the only method (to the best of our knowledge) to do so contrary to the other cell-cell interaction methods, which infer ligand-receptor interactions by employing the spatial data, as competitor methods, we selected state-of-the-art ligand-receptor interaction inference methods COMMOT \citep{Cang2023}, stLearn \citep{Pham2023} and SpatialDM \citep{Li2023}. These three methods are capable of quantifying ligand-receptor interactions at the resolution of spatial spots and we adopted them for quantifying ligand-target activity from our simulated datasets. The performance of each method was evaluated based on a spatial activity accuracy metric that computes the Pearson correlation of the spot-level ligand-target activity score measured by a method with the true simulated activity of the ligand-target pair in the positive and negative reference spatial locations (see Method for details). In addition, we also compared the methods in terms of precision and recall (see Methods) for the ligand-target pairs to evaluate their ability to distinguish the spots harboring ligand-target interactions from the spots devoid of such interactions.\par 

In each simulated dataset across the tissue types, Renoir was able to capture the ligand-target interactions more precisely. To illustrate, we consider the ligand-target pair \textit{A2M-MYC}, simulated from the human intestine dataset (Fig. 2b), where the positive (blue) and negative (red) reference spots included enrichment and abatement of said ligand-target pair interaction respectively as is also evident from the spatial expression plot of the two genes. Amongst the positive spots, we observed Renoir's ligand-target activity scores to truly reflect the enriched spots with high scores as compared to COMMOT and stLearn. While SpatialDM had high scores for positive spots, it also had scores all across the tissue and failed to differentiate the positive and negative spots appropriately. Renoir was also able to correctly quantify the absence of ligand-target interactions in the negative reference spots. In comparison, all three competitor methods showed false-positive activity of the gene pair in multiple negative spots (red bounded regions labeled 1, 2 and 3) with near obsolete \textit{A2M} and \textit{MYC} expression. Renoir is capable of computing the activity of a ligand-target by considering the expression of a suitable receptor at the receiving cell. To showcase this, we considered the pair \textit{IFNG-JCHAIN}, simulated using the TNBC dataset, where \textit{IFNG} was associated with Natural Killer T Cells (NKT Cells) and CD8$^{+}$ T Cells; and \textit{JCHAIN} was associated with B Cells and Dendritic Cells (DCs). However, B cells do not express a receptor for \textit{IFNG} (Supplementary Fig 1a) and thus the \textit{IFNG} secreted from NKT Cells or CD8$^{+}$ T Cells cannot activate \textit{JCHAIN} in B cells. Therefore, the positive reference spots harbored interactions with DCs as the target cell type whereas the negative spots contained the enrichment of B cells as the target cell type (see Methods section for more details). In this scenario, by accounting for the presence of suitable receptors while calculating ligand-target activity scores, Renoir correctly scored the negative spots with negligible scores; whereas other methods failed to consider such scenarios and therefore incorrectly gave high scores to the negative spots as depicted by the red bounded regions (Supplementary Fig 1a). Similarly, for the DLPFC dataset (e.g., for the pair \textit{NPY:CARTPT}) also, Renoir's spatial scores for the ligand-target activity were more precise as compared to the other methods (Supplementary Fig 1b). SpatialDM and COMMOT had less distinction between the positive and non-positive spots and stLearn failed to capture the activity pattern due to the high range of \textit{NPY} expression and relatively low range of \textit{CARTPT} expression. Quantitatively, Renoir outperformed other CCC inference methods in inferring spatial ligand-target activity by achieving 45.8, 98.5, and 12.4\% improvement in the median spatial activity accuracy across the human intestine, TNBC and DLPFC datasets (Fig 2c) indicating Renoir's superiority in capturing the spatial variation of the ligand-target activity. Furthermore, in terms of average precision and recall over the ligand-target pairs, Renoir outperformed SpatialDM, COMMOT, and stLearn across the tissue types across different threshold values used for calculating precision and recall (Supplementary Fig 1c). For varying threshold values, Renoir was able to consistently maintain high precision and recall values across the tissue types. In comparison, the precision and recall values for the other methods varied across datasets, with a method having high precision-recall for one dataset but low precision-recall for another. \par 

Finally, we evaluated the accuracy of the spatial communications domains inferred based on the ligand-target interaction scores inferred by a CCC inference method. For this, we utilized 10X Visium datasets from 12 dorsolateral prefrontal cortex (DLPFC) samples \citep{Maynard2021}, in which the authors manually annotated six cortical layers and white matter for each sample based on cytoarchitecture and selected marker genes. Considering the manual annotations as ground truth, we evaluated the utility of the ligand-target scores inferred by Renoir in identifying the layers and compared its performance against that of COMMOT, SpatialDM and stLearn. Each method was evaluated based on adjusted Rand index (ARI) and normalized mutual information (NMI) scores calculated by comparing the communication domains inferred by the method against the ground truth annotations over all 12 DLPFC samples (see Methods for details). Renoir achieved the highest ARI and NMI scores considering all the samples and outperformed all the other methods (Fig. 2d) indicating that it is able to capture spatial variation and architecture better than the other CCC inference methods. Figure 2e shows how the communication domains inferred by Renoir captured the underlying spatial structure for the sample 151675 better than the other methods for which the inferred communication domains were much noisier.

%% Renoir applied on mouse brain visium
\subsection*{Renoir identifies distinct communication domains in Mouse brain}

To demonstrate Renoir's ability to uncover region-specific ligand-target interactions, we first applied Renoir to 10X Visium data from two Mouse brain sections (samples S1 and S2) and its corresponding single nuclei data \citep{Kleshchevnikov2022} (see data availability). Renoir identified $9409$ (S1) and $9405$ (S2) ligand-target pairs that were expressed in this dataset and quantified their spatial neighborhood activity scores for both samples. Based on the ligand-target neighborhood activity scores, Renoir clustered the ST spots across the two tissue sections and identified $19$ spatial communication domains which had distinct latent space embeddings (Fig 3a) and were similarly distributed across the sections (Fig 3b). We further quantified the distribution of the spatial communication domains across different anatomical brain regions and found that the domains across both samples primarily corresponded to the same anatomical regions and were highly concordant across the two samples in terms of their distribution in anatomical brain regions (Fig 3c).\par

The communication domains were found to harbor diverse cell types that were spatially localized and abundance pattern concordant across the replicates (Supplementary Fig 2a-b). The domains further harbored differential activity of specific ligand-target pairs (Fig 2d, Supplementary Fig 3b-c). For instance, domain 1 specific to white matter exhibited \textit{Il33} and \textit{Csf1} activity which are known to be required for microglia development and maintenance in white matter \citep{10.3389/fimmu.2019.02199, 10.3389/fimmu.2018.02596} and domain 3 specific to the hypothalamus exhibited \textit{Dlk1} activity known to be expressed in hypothalamic arginine vasopressin and oxytocin neurons \citep{10.1371/journal.pone.0036134}.  We also noticed the existence of ligand or target-specific activity pertaining to distinct anatomical regions, wherein the ligand or target was known to be abundant within the brain regions corresponding to the domain. For example, domain 9 exhibited high activity of \textit{Rspo1}, which is known to be highly expressed within the cortex \citep{10.1210/en.2013-1550, Miller2019}, domain 7 exhibited high activity of hippocampal marker \textit{Crlf1}, and domain 13 exhibited high activity of \textit{Ccnd1} and \textit{Nr2f2}, which are known to be highly expressed within the Amygdala region \citep{KOELLER2008230, 10.1242/dev.075564}. Supplementary Table 1 lists the domain-specific ligand-target interactions with literature support for the associated brain region. \par 

% Domains 2 and \textcolor{blue}{10} located in the thalamus region were enriched in \textit{Tcf7l2} and \textit{Lef1} activity which are known to be essential for neurogenesis in the developing mouse neocortex \citep{Chodelkova2018, Nagalski2013}. \textcolor{blue}{domain 6 exhibited high \textit{Tgfa} activity, where TGF-$\alpha$ is synthesised across regions which included the tuber olfactorium, a part of the ventral striatum \citep{Wilcox1901},} 

Kleshchevnikov et al. \citep{Kleshchevnikov2022} identified and molecularly characterized 10 regional astrocyte subtypes using this spatial and matched single-cell data. However, interactions of these regional astrocyte subtypes with other spatially co-localized cell types were not explored. We applied Renoir to characterize the interaction of these regional astrocyte subtypes with other cell types. Figure 3d and Supplementary Figure 3d illustrate the interactions where ligands expressed by regional astrocyte subtypes activate other target genes in other cell types across samples S1 and S2 respectively. Figure 3e and Supplementary Figure 3e demonstrate the interactions where ligands expressed by other cell types activate specific marker genes in astrocyte subtypes. Ligands expressed by the telencephalic astrocytes in Cortex, and Amygdala were found to activate genes in excitatory neurons in Layer 6b (\textit{Vegfa:Ctgf}) and Amygdala (\textit{Ptn:Nr2f2}) respectively. On the other hand, \textit{Agt} expressed by diencephalic astrocytes (Thalamus) were found to activate genes in inhibitory neurons (\textit{Psap:Pde4dip}). White matter astrocytes interacted with oligodendrocytes via \textit{Tnc} which activated gene \textit{Adamts4}. Similarly, different genes in telencephalic astrocytes in Amygdala, Cortex, and Hippocampus were found to be activated by ligands expressed by excitatory neurons in Layer \textit{4} (\textit{Rspo1:Vegfa}), Amygdala (\textit{Slit1:Nr2f2}), and Hippocampus DG (\textit{Crlf1:Ccnd3}, \textit{Slit1:Tcf4}) respectively. Excitatory neurons in the thalamus were found to activate genes in diencephalic astrocytes (\textit{Ntng1:Gpld1}, \textit{Nxph1:Tcf7l2}).

For each communication domain, based on the agreement of neighborhood activity and prior regulatory potential scores, we ranked the prominent ligands (Fig 3f, Supplementary Fig 3f) that are expressed by the major cell types in the domain (see Methods for details) which further highlighted the region-specific enrichment of certain ligands. For instance, \textit{Bdnf} known to be associated with numerous functions within the hippocampus \citep{Heldt2007} displayed high activity within domain 16 associated with the hippocampus. \textit{Nrg1}, known to have high impact on cortical functions \citep{AGARWAL20141130} exhibited high activity within domain 5. Supplementary Table 2 lists such major ligands with their corresponding literature support. Using Renoir, we spatially mapped major pathway activities based on the activity scores of the ligand-target pairs associated with the pathways (see Methods) (Fig 3g, Supplementary Fig 3g and Supplementary Fig 4). Consistent with prior works \citep{https://doi.org/10.1002/cne.21186,10.3389/fnins.2016.00026}, cortex and hippocampus were highly enriched in MAPK signalling and NEUROTROPHIN signalling pathways, respectively. 

% This is consistent with prior works which demonstrated high concentrations of MAPK components in the cortex and hippocampus of the adult mouse brain , and NEUROTROPHIN signaling's role in the regulation of neurogenesis \citep{10.3389/fnins.2016.00026}. 

%% Renoir applied on TNBC visium
\subsection*{Renoir uncovers intratumor spatial communication sub-domains within triple negative breast cancer}

To demonstrate Renoir's ability to characterize complex heterogeneous tissues in terms of ligand-target interactions, we analyzed single-cell and Visium datasets (Fig 4a) from a triple negative breast cancer (TNBC) patient \citep{Wu2021}. Renoir quantified the spatial neighborhood activity scores for $18654$ ligand-target pairs based on which, it clustered the spatial spots into six communication domains which in addition to being distinct in low-dimensional embeddings (Fig 3b), were primarily associated with morphologically distinct pathological regions (Fig 3c). The domains were further differentiated based on the ligand-target activities and the composition and abundance of the cell types they harbored (Fig 4d-e, Supplementary Fig 5a-b). Since two (clusters 4 and 5) of the six domains had low overall UMI count (Supplementary Fig 5b), we focused on the first four communication domains in our further analyses. In accordance with their respective cell type compositions, interesting differential ligand-target activities were observed across the six domains. For instance, domain 0, histologically annotated as ductal carcinoma in situ (DCIS), harbored a high abundance of cancer cycling and cancer basal SC cell types (Fig 4c, Supplementary Fig 5a); cancer-associated fibroblasts (CAFs) and different T cell populations (Supplementary Fig 5a); and demonstrated differential activity of MYC and SAA1 (\textit{LAMB1-MYC}, \textit{SERPINE2-MYC}, \textit{SAA1-ESRRA}) known to be overexpressed in basal-subtype of breast cancer \citep{doi:10.1073/pnas.191367098,doi:10.1073/pnas.1732912100,OT26566}; and MDK (\textit{MDK-CCNA2}), a known cancer progression marker \citep{Filippou2020}. Domain 0 also harbored high activity of \textit{VEGF} (\textit{VEGFA}, \textit{VEGFB}) (known to be highly up-regulated in breast cancer \citep{RIBATTI2016453,Harmey1998,Lewis2000}), and other stromal ligands (\textit{COL18A1}, \textit{COL4A1}) (Fig 4e). Domain 1 was associated with the histological regions lymphocytes and invasive cancer and majorly harbored cancer-associated myofibroblasts (myCAF), macrophages, and T cells (cycling and \textit{CD8}$^{+}$) along with the two cancer cell types (Fig 4c, Supplementary Fig 5a). Consequently, interactions associated with tumor progression and invasion (e.g., \textit{CXCL10-MMP2}, \textit{CCL3-ICAM1} \citep{doi:10.1073/pnas.1408556111}, \textit{FN1-MMP10} \citep{MCGOWAN20081566}), as well as suppression (e.g., \textit{CXCL9-IRF1} \citep{10.1158/1078-0432.CCR-19-1868}) among these cell types were enriched in domain 1 (Fig 4e). Domain 2 was associated with the histological region normal, stroma and lymphocytes and harbored four major cell types including inflammatory-like CAFs (iCAF-like) and myoepithelial cells (Fig 4c, Supplementary Fig 5a). Interestingly, this domain harbored high activity of the ligand \textit{CCL28} which is known to promote breast cancer and metastasis \citep{Yang2017}, target \textit{HES1} which plays a crucial role in malignant cell proliferation \citep{Li2018} and \textit{AREG} which is expressed by myoepithelial cells and loss of which can inhibit tumor proliferation, growth and invasiveness \citep{Mao2018} (Fig 4e). Domains 3 and 4, despite low UMI count, showcased distinct ligand-target activities (e.g., \textit{IL6}, \textit{IL33}, \textit{IL1RN} in domain 3 and \textit{PENK} in domain 4 (Fig 4e)).\par 

We further investigated the most prominent cell type-specific ligand-target interactions (Fig 4f, Supplementary Fig 6a) within each spatial communication domain (see Methods for details). By emphasizing domains with the highest ligand-target activities and UMI count (0, 1, 2, and 3), we uncovered cell type-specific interactions such as \textit{RARRES2-AEBP1}, \textit{CXCL12-GSN} (ligand and target are markers of iCAFs, domain 2), \textit{CXCL16-CCND1} (domain 3). For each communication domain, we further ranked the major ligands (expressed by the major cell types in the domain) according to their activities (see Methods for details). A key EMT marker, \textit{TGFB3} \citep{jcm5070063} was one of the top-ranked ligands in both domains 0 and 1 (which harbored a high abundance of cancer cells) and its top targets included \textit{FN1} (in myCAFs and macrophages) and \textit{TNC} (in myoepithelial cells and myCAFs) which are associated with breast cancer metastasis \citep{Wang2018,Oskarsson2011}, (Fig 4g). In domain 1 which contained high abundance of macrophages, another top ligand was \textit{CCL2}, which is known to promote metastasis via retention of metastasis-associated macrophages (MAMs) \citep{10.1084/jem.20141836} and it was found to target \textit{ICAM1} (a molecular target for TNBC) \citep{doi:10.1073/pnas.1408556111} in myoepithelial cells and macrophages, other chemokines \textit{CCL3} and \textit{CCL4} as well as lymphocyte inhibitory molecule \textit{LTB} in \textit{CD4}$^{+}$ T cells (Fig 4g). Within each communication domain, we also found ligands that were specifically active in a single communication domain (Fig 4h). For example, in domain 1, \textit{CD8}$^{+}$ T cells secreted the ligand \textit{CCL5} (known to promote TNBC progression through the regulation of immunosuppressive cells \citep{Zhang2013}) which activated \textit{MMP9} and \textit{CCL2} in macrophages. \par

% \textcolor{blue}{We also noticed cross-communication between myCAFs and \textit{CD4}$^{+}$ T cells (\textit{CXCL12-FASLG}) in domain 1 [it was FASLG instead of ITK]}. Tumor metastasis and progression-promoting interactions between macrophages and cancer cycling cells (\textit{CXCL16-CCND1}) \citep{OT3690, cancers13143568} were observed in domain 3. Additionally in domain 3, we observed interaction between \textit{TGFB1} and \textit{FOXP3} which are markers of NK cells and \textit{CD4}$^{+}$ T cells respectively indicating \textit{TGFB1}-mediated suppression of immune response \citep{10.1371/journal.pone.0076379, doi.org/10.1111/imm.12941}. \par   

 % \textcolor{blue}{[I removed the last line here that emphasized on how CCL5 was distinct in domain 1]}

Subsequently, utilising cell type abundance values of spatial locations, we sub-clustered each communication domain (domains 0, 1, 2, and 3) into subdomains that were unique in terms of cell type abundance and entailed distinct ligand-target activities (Supplementary Figs 7-10). Amongst the six subdomains in domain 0 (Supplementary Fig 7), in subdomain 4 \textit{TGFB1} and \textit{FBN1} activated several EMT-associated genes (\textit{FBN1}, \textit{MMP2}, \textit{COL5A1}, \textit{PDGFRB}, etc.) \citep{10.1172/JCI45014} and cancer progression and regulation markers (\textit{DUSP1}, \textit{JUNB}, and \textit{MYC}) \citep{Tuglu2018, Lukey2016} respectively. In contrast, subdomain 1 which harbored different immune cells (\textit{CD4}$^{+}$ and \textit{CD8}$^{+}$ T cells, DCs, NKT cells) as well as iCAFs, showcased extensive chemokine activity (\textit{CCL5}) which targeted different genes in different immune cells and iCAFs (e.g., \textit{EGR1}). Domain 1 was segregated into 4 subdomains (Supplementary Fig 8), of which subdomain 2 abundant with immune cells showcased high \textit{RARRES2} activity (secreted by myCAFs) and subsequently a lower abundance of cancer cells compared to other subdomains. This is coherent with previous findings which show that \textit{RARRES2} serves as an immune-dependent tumor suppressor by acting as a chemo-attractant to recruit anti-cancer immune cells \citep{wittamer2003specific,wittamer2005neutrophil,zabel2005chemokine,vermi2005role}. Similarly, subdomains of other domains showcased interesting activity particular to the cell type abundant across the subdomains (Supplementary Figs 9-10). 
% Subdomain \textcolor{blue}{2} of domain 2 harbored a high abundance of macrophages, myoepithelial cells, myCAFs, and cycling myeloid cells displayed \textit{FBN1} and \textit{IL1B}-mediated activation of \textit{MMP1}, known to promote brain metastasis in TNBC \citep{Liu2012} and Domain 3 (segregated into five subdomains (Supplementary Figure 12)) of which subdomain \textcolor{blue}{4} abundant with immune cells showed high \textit{CCL19} \textcolor{blue}{\sout{and \textit{CCL21}}} activity which are known to attract lymphocytes \citep{vicari2000antitumor} to name a few.

Using Renoir, we further spatially mapped the major pathway activities based on the activity scores of the ligand-target pairs associated with the pathways (see Methods). While some pathways showed activity across multiple communication domains, some other pathways were more active in specific communication domains (Supplementary Fig 6b). For example, the WNT signaling pathway was found to be active in subdomain 1 of domain 0, and subdomain 2 of domain 1 which harbored a high abundance of cancer cell types (Fig 4i). IL-6/JAK/STAT3 signaling was highly active in subdomain 1 of domain 1 (Fig 4i) that harbored a high abundance of myCAFs, macrophages and myoepithelial cells along with cancer cells but lacked immune cells other than cycling T-cells indicating a highly immunosuppressive tumour microenvironment \citep{Johnson2018}. Taken together, Renoir was able to provide comprehensive insights into the spatial heterogeneity of ligand-target activity and its implication in tumor biology.

% TGF beta signaling was active in domains 0 and 1, and MAPK signaling had the highest activity in domain 2 (Figure 4i). 

\subsection*{Renoir identifies hepatocyte-macrophage interactions in developing fetal liver}

Cell-cell communication shapes tissue functioning and organogenesis during development. In order to evaluate the versatility of Renoir in uncovering cell-cell interactions during development, we applied Renoir to data from fetal liver development. We generated a 10X Visium spatial transcriptomic dataset (see Methods for details) from a fetal liver sample for which single-cell RNA sequencing data was available \citep{SHARMA2020377} (Fig 5a). Renoir identified $17630$ ligand-target pairs active in the tissue, quantified their spatial neighborhood activity scores and based on these scores, identified five spatial communication domains (Fig 5b), which differed based on ligand-target activity and distribution of UMI count across the tissue section (Fig 5b-c and Supplementary Fig 4a). Domain 0 demonstrated the majority of ligand-target activity, followed by domain 1. On the other hand, domains 2 and 4 showed sparse activity (UMI counts were low) (Fig 5c) while domain 3 was specific to the target \textit{CDKN1A} (Supplementary Fig 4a). This is in line with the UMI count distribution across spots, where the UMI count was most abundant in domains 0 and 3, followed by domains 1, 2 and 4. Cell types were found to be evenly distributed across domains 0, 1 and 3 (Supplementary Fig 11b). Activities pertinent to fetal liver were found to be widespread (Fig 5d). For example, \textit{COL18A1} associated with type XVIII collagen (a prominent component of the Extracellular Matrix (ECM) in the liver \citep{10.1242/dmm.011577}) was found to be highly expressed across the tissue and activate multiple targets including hematopoietic transcription factors \textit{BCL2L1} and \textit{BCL6} \citep{10.1371/journal.pone.0007641}. Higher \textit{CDKN1A} activity was observed, in line with the proliferative nature of hepatocytes in developing fetal liver tissue.\par
Further analyses involved acquiring previously described ligand rankings and prominent cell type-specific ligand-target interactions for each communication domain. Specifically, we observed prominent activity of the ligands \textit{MDK}, an important growth factor for fetal liver hematopoietic stem cells and multipotent progenitor expansion \citep{Gao2022} and \textit{CXCL12} which helps in retaining hematopoietic stem cells in the fetal liver \citep{Lewis2021} (Fig 5e). Importantly, we observed \textit{MARCO}$^+$ tissue-resident macrophage population (Kupffer cells) interacting with \textit{APOB}$^+$ hepatocytes. \textit{MARCO} (MAcrophage Receptor with COllagenous structure) is predominantly expressed in non-inflammatory Kupffer cells \citep{MacParland2018} and fetal-like reprogramming of liver macrophages plays an important role in immunosuppressive phenotype in liver cancer \citep{SHARMA2020377}. As shown in Fig 5e and Supplementary Fig 12a, we observed that \textit{MARCO}$^+$ macrophages interact with \textit{PLG}$^+$/\textit{APOB}$^+$ subset of hepatocytes. The \textit{PLG} (plasminogen) has been known to be important for hepatic regeneration \citep{doi:10.1073/pnas.96.26.15143,00001721-200809000-00006} and \textit{PLG} deficiency can lead to impaired remodelling in mouse models after acute injury. Next, we subdivided hepatocytes into two groups: \textit{ALB}$^+$\textit{APOA}$^+$\textit{PLG}$^+$ and \textit{ALB}$^+$\textit{APOA}$^+$\textit{PLG}$^-$ (Figs 4f-g). As shown in Figs 5g-h, these two subsets are transcriptionally and spatially distinct from each other indicating heterogeneity of the fetal liver hepatocyte population. We observed that the majority of the \textit{ALB}$^+$\textit{APOA}$^+$ hepatocyte population does not express \textit{PLG} while there is a specific cluster with a high expression of \textit{PLG} (Fig 5f). In the Visium data, the spots expressing \textit{ALB}$^+$\textit{APOA}$^+$\textit{PLG}$^+$ hepatocytes also contained \textit{FOLR2}$^+$ macrophages (Fig 5h). \textit{FOLR2}$^+$ macrophages are embryonic macrophages present in the fetal liver microenvironment \citep{SHARMA2020377}. Spatial correlation analyses further showed higher similarity of \textit{FOLR2}$^+$ macrophages and \textit{ALB}$^+$\textit{APOA}$^+$\textit{PLG}$^+$ hepatocytes as compared to \textit{PLG}$^-$ hepatocytes in terms of cell type distribution as well as neighborhood activity score (Fig 5i). Taken together, these results provide important insights into the macrophage-hepatocyte niche in the developing human liver. \par
% Major pathway activities were also spatially mapped onto the Fetal Liver sample (Supplementary Figure 5c). Although domain specific activities were undetected, relevant activity notably, KRAS signaling (plays an important role in Fetal Liver hematopoiesis \citep{10.1002/stem.2355}) and canonical JAK STAT signaling (recognised to be integral in initiating a signalling cascade required for embryonic development, haemopoietic development and differentiation \citep{doi:10.3109/08977194.2012.660936}) were distinctly observed. 

% We further observed the occurrence of ligand-target pairs \textit{CXCL12-CXCR4} and \textit{HGF-COL1A2} as prominent interactions (Supplementary Figure 5a). While \textit{CXCL12} is known to interact with \textit{CXCR4} to control the regeneration of liver amongst other organs \citep{10.3389/fimmu.2020.02109}, \textit{CXCR4} is known to be expressed at the highest density by the lymphoid cells \citep{https://doi.org/10.1002/SICI}. Moreover, the deletion of one of \textit{CXCL12} or \textit{CXCR4} is known to cause defective B-cell lymphopoiesis \citep{10.3389/fimmu.2020.02109}. 

\subsection*{Renoir identifies onco-fetal interactions in hepatocellular carcinoma}
To evaluate Renoir's ability in handling single-cell resolution ST datasets, we applied it on a Nanostring CosMx dataset generated from a hepatocellular carcinoma (HCC) patient. The dataset consisted of 70166 cells distributed across 46 cell types and 959 genes. 

We recently \citep{SHARMA2020377,Li2024} discovered the presence of onco-fetal tumor microenvironment (TME) in hepatocellular carcinoma, where the \textit{FOLR2}$^{+}$ TAMs, \textit{PLVAP}$^{+}$ ECs and \textit{POSTN}$^{+}$ CAFs were identified as onco-fetal cell types that shared transcriptional programs with their fetal liver counterparts. We evaluated Renoir's performance in identifying the onco-fetal interactions from the Single molecular imaging (CosMx) dataset. To perform an unbiased analysis, we used Renoir to quantify the activity of the ligands expressed by the onco-fetal cell types with the other major cell types present in the neighborhoods of the onco-fetal cells using Renoir's ligand-ranking formulation (Fig 6a). Renoir was able to characterize the activity of several known core onco-fetal ligands (e.g., \textit{VEGFA}, \textit{TGFB3}) as well as ligands shared by two onco-fetal cell types (e.g., \textit{BMP7}, \textit{CXCL12}, \textit{TGFB1}, \textit{IFNG}) (\citep{Li2024}). Specifically, Renoir identified the activity of known onco-fetal interactions \textit{CXCL12}-\textit{CXCR4} and \textit{VEGFA}-\textit{KDR} (Fig 6c). CXCL12-CXCR4 signaling was particularly elevated in oncofetal niche. Renoir correctly identified the co-localization of these cells and quantified the activity of the pair in the neighborhood. Similarly, VEGFA was expressed by all three onco-fetal cell types and it interacted with \textit{KDR} in co-localized ECs and \textit{MUC6}$^{+}$ bipotent hepatocytes (Fig 6c). 

Renior discovered oncofetal niche mediated reprogramming of MUC6+ bipotent cells

Renoir not only validated known oncofetal interaction but also identified novel biological interactions between oncofetal niche and putative stem-like MUC6+ cells in liver cancer microenvironments. Renoir identified the \textit{MUC6}$^{+}$ bipotent hepatocytes to have the highest number of interactions with the onco-fetal cell types (more specifically with \textit{POSTN}$^{+}$ CAFs) (Fig 6b). We delved into interactions specific to oncofetal cells and \textit{MUC6}$^{+}$ bipotent cells by quantifying the cell type-specific activity of the relevant ligand-target pairs in the onco-fetal neighborhoods (Supplementary Fig 13b). We noticed the activation of MYC and TP53 in \textit{MUC6}$^{+}$ bipotent cells by multiple onco-fetal ligands including \textit{BMP2, TGFB3, CXCL12}. We utilized gprofiler \citep{gProfiler} to infer the  pathways in the \textit{MUC6}$^{+}$ bipotent cells targeted by the oncofetal ligands and identified various interleukin signaling pathways (e.g., IL4, IL13, IL10), cytokine signaling and prostaglandin signaling (Fig 6d) to be enriched. Since prostaglandin signaling is known to promote HCC progression \citep{prostaglandin_pathway}, we further looked into the activity of \textit{IL6-PTGS2}. High scores for \textit{IL6-PTGS2} was found in the tissue niche where \textit{MUC6}$^{+}$ bipotent cells were co-localized with the oncofetal cells indicating the potential role of onco-fetal cells in the activation of PTGS2 in the \textit{MUC6}$^{+}$ bipotent cells (Fig 6e). Importantly, these observation indicate that application of Renoir could lead to identification of novel biological insights in tumor microenvironment. 